% Options for packages loaded elsewhere
\PassOptionsToPackage{unicode}{hyperref}
\PassOptionsToPackage{hyphens}{url}
%
\documentclass[
]{article}
\usepackage{amsmath,amssymb}
\usepackage{iftex}
\ifPDFTeX
  \usepackage[T1]{fontenc}
  \usepackage[utf8]{inputenc}
  \usepackage{textcomp} % provide euro and other symbols
\else % if luatex or xetex
  \usepackage{unicode-math} % this also loads fontspec
  \defaultfontfeatures{Scale=MatchLowercase}
  \defaultfontfeatures[\rmfamily]{Ligatures=TeX,Scale=1}
\fi
\usepackage{lmodern}
\ifPDFTeX\else
  % xetex/luatex font selection
\fi
% Use upquote if available, for straight quotes in verbatim environments
\IfFileExists{upquote.sty}{\usepackage{upquote}}{}
\IfFileExists{microtype.sty}{% use microtype if available
  \usepackage[]{microtype}
  \UseMicrotypeSet[protrusion]{basicmath} % disable protrusion for tt fonts
}{}
\makeatletter
\@ifundefined{KOMAClassName}{% if non-KOMA class
  \IfFileExists{parskip.sty}{%
    \usepackage{parskip}
  }{% else
    \setlength{\parindent}{0pt}
    \setlength{\parskip}{6pt plus 2pt minus 1pt}}
}{% if KOMA class
  \KOMAoptions{parskip=half}}
\makeatother
\usepackage{xcolor}
\usepackage{color}
\usepackage{fancyvrb}
\newcommand{\VerbBar}{|}
\newcommand{\VERB}{\Verb[commandchars=\\\{\}]}
\DefineVerbatimEnvironment{Highlighting}{Verbatim}{commandchars=\\\{\}}
% Add ',fontsize=\small' for more characters per line
\newenvironment{Shaded}{}{}
\newcommand{\AlertTok}[1]{\textcolor[rgb]{1.00,0.00,0.00}{\textbf{#1}}}
\newcommand{\AnnotationTok}[1]{\textcolor[rgb]{0.38,0.63,0.69}{\textbf{\textit{#1}}}}
\newcommand{\AttributeTok}[1]{\textcolor[rgb]{0.49,0.56,0.16}{#1}}
\newcommand{\BaseNTok}[1]{\textcolor[rgb]{0.25,0.63,0.44}{#1}}
\newcommand{\BuiltInTok}[1]{\textcolor[rgb]{0.00,0.50,0.00}{#1}}
\newcommand{\CharTok}[1]{\textcolor[rgb]{0.25,0.44,0.63}{#1}}
\newcommand{\CommentTok}[1]{\textcolor[rgb]{0.38,0.63,0.69}{\textit{#1}}}
\newcommand{\CommentVarTok}[1]{\textcolor[rgb]{0.38,0.63,0.69}{\textbf{\textit{#1}}}}
\newcommand{\ConstantTok}[1]{\textcolor[rgb]{0.53,0.00,0.00}{#1}}
\newcommand{\ControlFlowTok}[1]{\textcolor[rgb]{0.00,0.44,0.13}{\textbf{#1}}}
\newcommand{\DataTypeTok}[1]{\textcolor[rgb]{0.56,0.13,0.00}{#1}}
\newcommand{\DecValTok}[1]{\textcolor[rgb]{0.25,0.63,0.44}{#1}}
\newcommand{\DocumentationTok}[1]{\textcolor[rgb]{0.73,0.13,0.13}{\textit{#1}}}
\newcommand{\ErrorTok}[1]{\textcolor[rgb]{1.00,0.00,0.00}{\textbf{#1}}}
\newcommand{\ExtensionTok}[1]{#1}
\newcommand{\FloatTok}[1]{\textcolor[rgb]{0.25,0.63,0.44}{#1}}
\newcommand{\FunctionTok}[1]{\textcolor[rgb]{0.02,0.16,0.49}{#1}}
\newcommand{\ImportTok}[1]{\textcolor[rgb]{0.00,0.50,0.00}{\textbf{#1}}}
\newcommand{\InformationTok}[1]{\textcolor[rgb]{0.38,0.63,0.69}{\textbf{\textit{#1}}}}
\newcommand{\KeywordTok}[1]{\textcolor[rgb]{0.00,0.44,0.13}{\textbf{#1}}}
\newcommand{\NormalTok}[1]{#1}
\newcommand{\OperatorTok}[1]{\textcolor[rgb]{0.40,0.40,0.40}{#1}}
\newcommand{\OtherTok}[1]{\textcolor[rgb]{0.00,0.44,0.13}{#1}}
\newcommand{\PreprocessorTok}[1]{\textcolor[rgb]{0.74,0.48,0.00}{#1}}
\newcommand{\RegionMarkerTok}[1]{#1}
\newcommand{\SpecialCharTok}[1]{\textcolor[rgb]{0.25,0.44,0.63}{#1}}
\newcommand{\SpecialStringTok}[1]{\textcolor[rgb]{0.73,0.40,0.53}{#1}}
\newcommand{\StringTok}[1]{\textcolor[rgb]{0.25,0.44,0.63}{#1}}
\newcommand{\VariableTok}[1]{\textcolor[rgb]{0.10,0.09,0.49}{#1}}
\newcommand{\VerbatimStringTok}[1]{\textcolor[rgb]{0.25,0.44,0.63}{#1}}
\newcommand{\WarningTok}[1]{\textcolor[rgb]{0.38,0.63,0.69}{\textbf{\textit{#1}}}}
\setlength{\emergencystretch}{3em} % prevent overfull lines
\providecommand{\tightlist}{%
  \setlength{\itemsep}{0pt}\setlength{\parskip}{0pt}}
\setcounter{secnumdepth}{-\maxdimen} % remove section numbering
\ifLuaTeX
  \usepackage{selnolig}  % disable illegal ligatures
\fi
\IfFileExists{bookmark.sty}{\usepackage{bookmark}}{\usepackage{hyperref}}
\IfFileExists{xurl.sty}{\usepackage{xurl}}{} % add URL line breaks if available
\urlstyle{same}
\hypersetup{
  hidelinks,
  pdfcreator={LaTeX via pandoc}}

\author{}
\date{}

\begin{document}

\hypertarget{arquitectura-de-satspath-en-arkade-wallet}{%
\section{Arquitectura de SatsPath en Arkade
Wallet}\label{arquitectura-de-satspath-en-arkade-wallet}}

El protocolo de \textbf{SatsPath} está completamente integrado en la
wallet a través de un puente con WebAssembly (WASM). Esto permite que
toda la criptografía pesada y la lógica de ruteo y resolución ocurran de
forma ultra-rápida y segura del lado del cliente, utilizando el mismo
código robusto en Rust de \texttt{satspath-core}.

A continuación te presento un diagrama de secuencia de cómo fluyen los
datos en las dos etapas principales: el \textbf{Flujo del Receptor}
(creación de identidad) y el \textbf{Flujo del Emisor} (resolución y
envío).

\begin{Shaded}
\begin{Highlighting}[]
\NormalTok{sequenceDiagram}
\NormalTok{    participant User as Usuario}
\NormalTok{    participant Wallet as Arkade Wallet (React)}
\NormalTok{    participant Bridge as satspath.ts (WASM Bridge)}
\NormalTok{    participant WASM as satspath{-}wasm (Rust)}
\NormalTok{    participant Network as Red (Nostr / DNS)}
    
\NormalTok{    \%\% Flujo del Receptor}
\NormalTok{    rect rgb(35, 20, 50)}
\NormalTok{    Note over User, Network: 1. FLUJO DEL RECEPTOR (Crear Identidad)}
\NormalTok{    User{-}\textgreater{}\textgreater{}Wallet: Escribe Alias y Direcciones (Settings)}
\NormalTok{    Wallet{-}\textgreater{}\textgreater{}Wallet: Detona animación de Láseres 💥}
\NormalTok{    Wallet{-}\textgreater{}\textgreater{}Bridge: buildSatsPathProfile()}
\NormalTok{    Bridge{-}\textgreater{}\textgreater{}WASM: generate\_identity\_keypair()}
\NormalTok{    Note right of WASM: Genera llave secp256k1}
\NormalTok{    WASM{-}{-}\textgreater{}\textgreater{}Bridge: IdentityKeypair (Hex)}
\NormalTok{    Bridge{-}\textgreater{}\textgreater{}Bridge: Construye Perfil JSON crudo}
\NormalTok{    Bridge{-}\textgreater{}\textgreater{}WASM: sign\_profile\_json(json, secretKey)}
\NormalTok{    Note right of WASM: Aplica hash y firma de Schnorr}
\NormalTok{    WASM{-}{-}\textgreater{}\textgreater{}Bridge: Firma Criptográfica}
\NormalTok{    Bridge{-}{-}\textgreater{}\textgreater{}Wallet: SignedPaymentProfile}
\NormalTok{    Wallet{-}\textgreater{}\textgreater{}Wallet: Guarda en LocalStorage}
\NormalTok{    Wallet{-}\textgreater{}\textgreater{}User: Muestra JSON firmado listo para la red}
\NormalTok{    end}

\NormalTok{    \%\% Flujo del Emisor}
\NormalTok{    rect rgb(50, 30, 20)}
\NormalTok{    Note over User, Network: 2. FLUJO DEL EMISOR (Enviar Sats)}
\NormalTok{    User{-}\textgreater{}\textgreater{}Wallet: Enviar a "chelo@dev.idk"}
\NormalTok{    Wallet{-}\textgreater{}\textgreater{}Bridge: routePayment()}
\NormalTok{    Bridge{-}\textgreater{}\textgreater{}WASM: quote()}
\NormalTok{    WASM{-}\textgreater{}\textgreater{}Network: Resuelve Perfil (BIP353/NIP05/Local)}
\NormalTok{    Network{-}{-}\textgreater{}\textgreater{}WASM: SignedPaymentProfile}
\NormalTok{    Note right of WASM: Valida firmas y expiración}
\NormalTok{    WASM{-}\textgreater{}\textgreater{}WASM: router::select\_route()}
\NormalTok{    Note right of WASM: Sistema de Scoring elige \textless{}br/\textgreater{}el método más barato/óptimo}
\NormalTok{    WASM{-}{-}\textgreater{}\textgreater{}Bridge: QuoteResponse (Mejor método)}
\NormalTok{    Bridge{-}{-}\textgreater{}\textgreater{}Wallet: RouteResult (Ej. Lightning)}
\NormalTok{    Wallet{-}\textgreater{}\textgreater{}User: Muestra pantalla de confirmación}
\NormalTok{    end}
\end{Highlighting}
\end{Shaded}

\hypertarget{componentes-clave}{%
\subsection{Componentes Clave}\label{componentes-clave}}

\begin{enumerate}
\def\labelenumi{\arabic{enumi}.}
\tightlist
\item
  \textbf{Arkade Wallet (React):} Es la capa de presentación. Maneja las
  animaciones (como los láseres), la recolección de datos del usuario,
  el guardado local y muestra los resultados finales.
\item
  \textbf{\texttt{satspath.ts} (WASM Bridge):} Es el traductor. Se
  encarga de convertir los objetos de JavaScript a los tipos primitivos
  que espera Rust, y viceversa. Aquí es donde inyectamos los ``mocks''
  para desarrollo.
\item
  \textbf{\texttt{satspath-wasm} (Rust):} El cerebro de la operación.
  Hereda el código de \texttt{satspath-core}. Contiene el generador
  criptográfico seguro, la lógica matemática de las firmas de Schnorr, y
  el motor de Scoring que decide si es más barato irse por Ark,
  Lightning o On-chain dependiendo de las comisiones en tiempo real.
\end{enumerate}

\hypertarget{visiuxf3n-arquitectura-soberana-p2p}{%
\subsection{Visión: Arquitectura Soberana
P2P}\label{visiuxf3n-arquitectura-soberana-p2p}}

Este es el modelo teórico 100\% cypherpunk que busca eliminar la
dependencia de servidores DNS y HTTP, garantizando privacidad mediante
conexiones enjambre y Web of Trust.

\begin{Shaded}
\begin{Highlighting}[]
\NormalTok{graph TD}
\NormalTok{    \%\% Flujo de Usuarios}
\NormalTok{    User((Usuario)) {-}{-} "Enviar Sats" {-}{-}\textgreater{} Bob((Bob))}
\NormalTok{    Bob {-}{-} "Recibir Sats" {-}{-}\textgreater{} QR["Muestra Código QR"]}
\NormalTok{    QR {-}{-}\textgreater{} Alias["QR contiene Alias P2P"]}
\NormalTok{    Alias {-}{-}\textgreater{} Address["bob@satspath.p2p"]}
\NormalTok{    Address {-}{-}\textgreater{} PubKeys["Contiene Identity Pubkey"]}

\NormalTok{    \%\% Protocolos y Capas}
\NormalTok{    subgraph Capas de Red}
\NormalTok{        LN["Lightning Network"]}
\NormalTok{        ONCHAIN["On{-}chain"]}
\NormalTok{        ARK["Ark Protocol"]}
\NormalTok{        P2P["Red P2P (Hyperswarm)"]}
\NormalTok{    end}

\NormalTok{    \%\% Seguridad Soberana}
\NormalTok{    subgraph Security [Seguridad Soberana y Confianza P2P]}
\NormalTok{        Sec["Seguridad"]}
\NormalTok{        Enc["Cifrado P2P (ECDH)"]}
\NormalTok{        Id["Identidad Criptográfica (Pubkey secp256k1)"]}
        
\NormalTok{        \%\% Reemplazamos Email/Phone por métodos descentralizados}
\NormalTok{        QR\_Scan["Verificación Out{-}of{-}band (QR Scan)"]}
\NormalTok{        WoT["Web of Trust (Gráfica Social Nostr)"]}
\NormalTok{        Hashcash["Prueba de Trabajo (Anti{-}Spam)"]}
        
\NormalTok{        Trust["Confianza Matemática y Social"]}

\NormalTok{        Sec {-}{-}\textgreater{} Enc}
\NormalTok{        Sec {-}{-}\textgreater{} Id}
\NormalTok{        Sec {-}{-}\textgreater{} QR\_Scan}
\NormalTok{        Sec {-}{-}\textgreater{} WoT}
\NormalTok{        Sec {-}{-}\textgreater{} Hashcash}
        
\NormalTok{        QR\_Scan {-}{-}\textgreater{} Trust}
\NormalTok{        WoT {-}{-}\textgreater{} Trust}
\NormalTok{        Hashcash {-}{-}\textgreater{} Trust}
\NormalTok{    end}
\end{Highlighting}
\end{Shaded}

\begin{center}\rule{0.5\linewidth}{0.5pt}\end{center}

\hypertarget{arquitectura-actual-implementada-en-el-cuxf3digo}{%
\subsection{Arquitectura Actual (Implementada en el
Código)}\label{arquitectura-actual-implementada-en-el-cuxf3digo}}

A diferencia del diagrama teórico 100\% soberano que discutimos arriba,
el código actual de \texttt{satspath-core} y \texttt{satspath-router}
implementa una arquitectura híbrida de transición.

Actualmente, el sistema resuelve los alias apoyándose en
infraestructuras existentes (DNS y HTTP) y tiene un prototipo básico de
P2P con Pear/Hyperswarm, pero \textbf{aún no implementa} el cifrado
ECDH, los Handshakes ni las pruebas anti-spam (Hashcash).

Aquí tienes el diagrama exacto de cómo funciona el código \textbf{hoy}:

\begin{Shaded}
\begin{Highlighting}[]
\NormalTok{graph TD}
\NormalTok{    \%\% Flujo de Usuarios}
\NormalTok{    User((Usuario)) {-}{-} "Enviar Sats a chelo@dev.idk" {-}{-}\textgreater{} Wallet[(Arkade Wallet)]}
\NormalTok{    Wallet {-}{-} "routePayment()" {-}{-}\textgreater{} Bridge[satspath.ts WASM Bridge]}
\NormalTok{    Bridge {-}{-} "quote()" {-}{-}\textgreater{} Router\{"SatsPath Router (Rust)"\}}

\NormalTok{    \%\% Resolvers Actuales Implementados}
\NormalTok{    subgraph Resolvers [Módulos de Resolución Actuales]}
\NormalTok{        DNS["BIP353Resolver (DNS TXT)"]}
\NormalTok{        Nostr["NostrResolver (HTTP NIP{-}05)"]}
\NormalTok{        Pear["PearResolver (P2P Básico)"]}
\NormalTok{        Local["MockResolver (Desarrollo)"]}
\NormalTok{    end}

\NormalTok{    Router {-}{-}\textgreater{} DNS}
\NormalTok{    Router {-}{-}\textgreater{} Nostr}
\NormalTok{    Router {-}{-}\textgreater{} Pear}
\NormalTok{    Router {-}{-}\textgreater{} Local}

\NormalTok{    \%\% Detalle de los Resolvers}
\NormalTok{    DNS {-}. "Consulta DNS" .{-}\textgreater{} Cloudflare[Servidor DNS]}
\NormalTok{    Nostr {-}. "GET /.well{-}known/nostr.json" .{-}\textgreater{} WebServer[Servidor Web]}
    
\NormalTok{    \%\% Detalle del fallo de soberanía actual}
\NormalTok{    Pear {-}. "node satspath{-}pear/index.js resolve \textless{}alias\textgreater{}" .{-}\textgreater{} NodeJS[Script Local Node.js]}
\NormalTok{    NodeJS {-}. "SHA256(alias) = Topic" .{-}\textgreater{} Hyperswarm[Hyperswarm DHT]}
    
\NormalTok{    \%\% Proceso de Scoring}
\NormalTok{    Cloudflare {-}. "SignedPaymentProfile" .{-}\textgreater{} Router}
\NormalTok{    WebServer {-}. "SignedPaymentProfile" .{-}\textgreater{} Router}
\NormalTok{    Hyperswarm {-}. "SignedPaymentProfile" .{-}\textgreater{} Router}
    
\NormalTok{    Router {-}{-} "router::select\_route()" {-}{-}\textgreater{} Scoring[Motor de Scoring]}
\NormalTok{    Scoring {-}{-} "QuoteResponse" {-}{-}\textgreater{} Wallet}
\end{Highlighting}
\end{Shaded}

\begin{center}\rule{0.5\linewidth}{0.5pt}\end{center}

\hypertarget{implementation-plan-migraciuxf3n-a-p2p-soberano-trusted-connections}{%
\subsection{Implementation Plan: Migración a P2P Soberano (Trusted
Connections)}\label{implementation-plan-migraciuxf3n-a-p2p-soberano-trusted-connections}}

Para lograr que el código de SatsPath cumpla con el modelo 100\%
soberano (y resolver el problema de privacidad del alias expuesto), se
debe ejecutar el siguiente plan de desarrollo:

\hypertarget{fase-1-enjambres-ofuscados-y-cifrado-ecdh}{%
\subsubsection{Fase 1: Enjambres Ofuscados y Cifrado
(ECDH)}\label{fase-1-enjambres-ofuscados-y-cifrado-ecdh}}

\begin{enumerate}
\def\labelenumi{\arabic{enumi}.}
\tightlist
\item[$\square$]
  \textbf{Derivación de Tópico Seguro:} Modificar \texttt{satspath-core}
  para que no pase el alias en texto plano al script de Node.
\item[$\square$]
  \textbf{ECDH Compartido:} Implementar Diffie-Hellman usando la
  \texttt{IdentityPubkey} del emisor y receptor para generar un secreto.
  El \texttt{Topic} de Hyperswarm será el \texttt{SHA256(SharedSecret)}.
\item[$\square$]
  \textbf{Cifrado de Carga Útil:} Usar ChaCha20-Poly1305 (con el secreto
  compartido) para encriptar los paquetes de datos P2P, garantizando
  confidencialidad total.
\end{enumerate}

\hypertarget{fase-2-protocolo-tcp-style-syn-syn-ack}{%
\subsubsection{Fase 2: Protocolo TCP-Style (SYN /
SYN-ACK)}\label{fase-2-protocolo-tcp-style-syn-syn-ack}}

\begin{enumerate}
\def\labelenumi{\arabic{enumi}.}
\tightlist
\item[$\square$]
  \textbf{Paquete SYN (Emisor -\textgreater{} Receptor):} El nodo de
  Hyperswarm emisor debe enviar un payload cifrado con su intención de
  pago (\texttt{Amount} y \texttt{Nonce} aleatorio), y firmado
  criptográficamente por la Wallet.
\item[$\square$]
  \textbf{Generación Just-In-Time (Ark):} El nodo receptor recibe el
  SYN, valida la firma, y detona un llamado al servidor ASP de Ark para
  generar un VTXO dinámico que se bloquee específicamente para esta
  transacción.
\item[$\square$]
  \textbf{Paquete SYN-ACK (Receptor -\textgreater{} Emisor):} El
  receptor devuelve el invoice fresco, su alias local, y firma el
  \texttt{Nonce} para probar liveness.
\end{enumerate}

\hypertarget{fase-3-escudos-anti-bot-anti-ia-web-of-trust}{%
\subsubsection{Fase 3: Escudos Anti-Bot / Anti-IA (Web of
Trust)}\label{fase-3-escudos-anti-bot-anti-ia-web-of-trust}}

\begin{enumerate}
\def\labelenumi{\arabic{enumi}.}
\tightlist
\item[$\square$]
  \textbf{Módulo de Web of Trust (WoT):} Integrar Nostr (NIP-65) en la
  wallet para mantener una gráfica social de llaves públicas conocidas.
\item[$\square$]
  \textbf{Reglas de Aceptación:} Configurar el demonio P2P para
  rechazar/dropear automáticamente conexiones \texttt{SYN} de Pubkeys
  que estén a más de 2 grados de separación en la red de confianza.
\item[$\square$]
  \textbf{Integración de Hashcash (Lightning/Ark):} Obligar a conexiones
  fuera de la WoT a adjuntar un HTLC válido por una micro-transacción de
  10 sats para cubrir el costo computacional del ataque Sybil.
\end{enumerate}

\end{document}
